\documentclass[12pt,a4paper]{article}
\usepackage[utf8]{inputenc}
\usepackage[vietnamese]{babel}
\usepackage{mathptmx} % Font Times New Roman
\usepackage[a4paper,top=2cm,bottom=2cm,left=2cm,right=2cm]{geometry}
\usepackage{amsmath,amssymb}
\usepackage{listings}
\usepackage{xcolor}
\usepackage{tikz}
\usetikzlibrary{shapes.geometric,shapes.multipart,arrows.meta,positioning,calc,fit}
\usepackage{booktabs}
\usepackage{tabularx}
\usepackage{array}
\usepackage{enumitem}
\usepackage{float}
\usepackage{needspace}
\usepackage{setspace}
\setstretch{1.14}

% Định nghĩa màu sắc (đồng bộ với hệ thống bài giảng)
\definecolor{codegreen}{rgb}{0,0.5,0}
\definecolor{codegray}{rgb}{0.4,0.4,0.4}
\definecolor{codepurple}{rgb}{0.55,0,0.85}
\definecolor{codeblue}{rgb}{0.0,0.22,0.65}
\definecolor{backcolour}{rgb}{0.97,0.97,0.98}
\definecolor{addrcolor}{rgb}{0.75,0.0,0.0}
\definecolor{stepblue}{rgb}{0.1,0.2,0.7}

% Thiết lập listings cho C++
\lstset{
    language=C++,
    backgroundcolor=\color{backcolour},
    commentstyle=\color{codegreen}\itshape,
    keywordstyle=\color{codeblue}\bfseries,
    numberstyle=\tiny\color{codegray},
    stringstyle=\color{codepurple},
    basicstyle=\fontfamily{pcr}\footnotesize,
    breakatwhitespace=false,
    breaklines=true,
    captionpos=b,
    keepspaces=true,
    numbers=left,
    numbersep=6pt,
    showspaces=false,
    showstringspaces=false,
    showtabs=false,
    tabsize=4,
    frame=single,
    rulecolor=\color{codegray!35},
    aboveskip=0.5em,
    belowskip=0.5em,
    escapeinside={(*@}{@*)}
}

% Định dạng tiêu đề các mục
\usepackage{titlesec}
\titleformat{\section}{\Large\bfseries\color{codeblue}}{\thesection.}{0.6em}{}
\titleformat{\subsection}{\large\bfseries}{\thesubsection.}{0.5em}{}
\titleformat{\subsubsection}{\normalsize\bfseries}{\thesubsubsection.}{0.5em}{}
\titlespacing*{\section}{0pt}{1.2em}{0.5em}
\titlespacing*{\subsection}{0pt}{1.0em}{0.4em}
\titlespacing*{\subsubsection}{0pt}{0.8em}{0.3em}

\begin{document}

\begin{center}
    {\large \textbf{BÀI GIẢNG HỌC PHẦN: CẤU TRÚC DỮ LIỆU VÀ GIẢI THUẬT}}\\[0.25em]
    {\LARGE \textbf{\color{codeblue}CHUYÊN ĐỀ: DANH SÁCH LIÊN KẾT ĐƠN}}\\[0.35em]
    \rule{\textwidth}{1pt}
\end{center}

\vspace{-0.6em}

\section*{TỔNG QUAN NỘI DUNG BÀI HỌC}
\begin{enumerate}[leftmargin=2.5em, itemsep=0.05em]
    \item Ý tưởng -- Giải pháp
    \item Kiến trúc danh sách liên kết đơn
    \item Các bước xây dựng
    \item Các thao tác
    \item Bài tập
\end{enumerate}

\vspace{0.4em}
\hrule
\vspace{0.8em}

% ==============================================================================
\section{Ý Tưởng -- Giải Pháp}
% ==============================================================================

\subsection{Vấn đề 1: Lưu trữ một số lượng nhỏ các giá trị}
\textbf{Bài toán đặt ra:} Muốn lưu trữ các số: \texttt{10}, \texttt{79}, \texttt{39}, \texttt{66}, \texttt{26}, \texttt{50} thì chúng ta làm sao?

\noindent Một cách tiếp cận đơn giản là khai báo các biến rời rạc tương ứng cho từng giá trị:
\begin{center}
\begin{tabular}{cccccc}
    \texttt{10} & \texttt{79} & \texttt{39} & \texttt{66} & \texttt{26} & \texttt{50} \\
    \texttt{x}  & \texttt{y}  & \texttt{z}  & \texttt{t}  & \texttt{g}  & \texttt{h}
\end{tabular}
\end{center}

\subsection{Vấn đề 2: Lưu trữ nhiều số hơn}
\textbf{Bài toán đặt ra:} Muốn lưu trữ nhiều số hơn nữa thì làm sao?

\noindent Giải pháp quen thuộc là sử dụng \textbf{mảng một chiều} (ví dụ cấp phát mảng gồm 500 phần tử có chỉ số từ \texttt{0} đến \texttt{499}):

\begin{table}[H]
\centering
\renewcommand{\arraystretch}{1.25}
\begin{tabular}{|>{\bfseries}l|c|c|c|c|c|c|c|c|c|}
\hline
\textbf{index} & 0 & 1 & 2 & 3 & 4 & 5 & 6 & \dots & 499 \\
\hline
\textbf{value} & 10 & 79 & 39 & 66 & 26 & 50 &  & \dots &  \\
\hline
\end{tabular}
\caption{Minh họa lưu trữ dãy số bằng mảng một chiều.}
\end{table}

\noindent\textbf{Câu hỏi mở:} \textit{Có giải pháp nào tốt hơn nữa không?} (Việc cấp phát mảng cố định kích thước có thể gây dư thừa vùng nhớ chưa dùng đến hoặc không đủ chỗ khi số lượng phần tử tăng thêm $\Rightarrow$ cần cấu trúc dữ liệu cấp phát linh hoạt theo từng phần tử là \textbf{Danh sách liên kết đơn}).

% ==============================================================================
\section{Kiến Trúc Danh Sách Liên Kết Đơn}
% ==============================================================================

Xét nhu cầu lưu trữ dãy các số nguyên: \texttt{10}, \texttt{79}, \texttt{39}, \texttt{66}, \texttt{26}.

Thay vì lưu các phần tử liên tiếp nhau trong bộ nhớ như mảng một chiều, mỗi phần tử trong \textbf{danh sách liên kết đơn} được đóng gói thành một \textbf{Node} gồm hai thành phần:
\begin{itemize}[itemsep=0.15em]
    \item \texttt{data}: Lưu trữ dữ liệu (giá trị) của phần tử.
    \item \texttt{pNext}: Con trỏ lưu địa chỉ của \texttt{Node} kế tiếp trong danh sách (Node cuối cùng trỏ tới \texttt{NULL}).
\end{itemize}
Toàn bộ danh sách được quản lý thông qua hai con trỏ đặc biệt:
\begin{itemize}[itemsep=0.15em]
    \item \texttt{pHead}: Con trỏ trỏ đến \texttt{Node} đầu tiên của danh sách.
    \item \texttt{pTail}: Con trỏ trỏ đến \texttt{Node} cuối cùng của danh sách.
\end{itemize}

\begin{figure}[H]
\centering
\begin{tikzpicture}[
    nbox/.style={rectangle, draw=blue!70!black, fill=white, minimum width=0.82cm, minimum height=0.5cm, font=\small, inner sep=2pt},
    pbox/.style={rectangle, draw=orange!80!black, fill=white, minimum width=0.82cm, minimum height=0.5cm, font=\footnotesize, inner sep=2pt},
    lbl/.style={font=\ttfamily\tiny, text=codeblue},
    addr/.style={font=\ttfamily\footnotesize, text=addrcolor},
    ptr/.style={font=\bfseries\small, text=addrcolor}
]
% Các giá trị gốc phía trên
\foreach \val/\x in {10/0, 79/3.0, 39/6.0, 66/9.0, 26/12.0} {
    \node[font=\bfseries\color{codeblue}] at (\x+0.42, 1.85) {\val};
}

% Node 1: 0xb12
\node[nbox] (d1) at (0, 0) {10};
\node[pbox, right=0.05cm of d1] (n1) {a78};
\node[lbl, anchor=south] at (d1.north) {data};
\node[lbl, anchor=south] at (n1.north) {pNext};
\node[ellipse, draw=red!80!black, thick, inner sep=3.5pt, fit=(d1)(n1), yshift=0.14cm] (b1) {};
\node[addr, anchor=south] at (b1.north) {0xb12};

% Node 2: 0xa78
\node[nbox] (d2) at (3.0, 0) {79};
\node[pbox, right=0.05cm of d2] (n2) {d56};
\node[lbl, anchor=south] at (d2.north) {data};
\node[lbl, anchor=south] at (n2.north) {pNext};
\node[ellipse, draw=red!80!black, thick, inner sep=3.5pt, fit=(d2)(n2), yshift=0.14cm] (b2) {};
\node[addr, anchor=south] at (b2.north) {0xa78};

% Node 3: 0xd56
\node[nbox] (d3) at (6.0, 0) {39};
\node[pbox, right=0.05cm of d3] (n3) {c95};
\node[lbl, anchor=south] at (d3.north) {data};
\node[lbl, anchor=south] at (n3.north) {pNext};
\node[ellipse, draw=red!80!black, thick, inner sep=3.5pt, fit=(d3)(n3), yshift=0.14cm] (b3) {};
\node[addr, anchor=south] at (b3.north) {0xd56};

% Node 4: 0xc95
\node[nbox] (d4) at (9.0, 0) {66};
\node[pbox, right=0.05cm of d4] (n4) {a36};
\node[lbl, anchor=south] at (d4.north) {data};
\node[lbl, anchor=south] at (n4.north) {pNext};
\node[ellipse, draw=red!80!black, thick, inner sep=3.5pt, fit=(d4)(n4), yshift=0.14cm] (b4) {};
\node[addr, anchor=south] at (b4.north) {0xc95};

% Node 5: 0xa36
\node[nbox] (d5) at (12.0, 0) {26};
\node[pbox, fill=yellow!80, right=0.05cm of d5] (n5) {NULL};
\node[lbl, anchor=south] at (d5.north) {data};
\node[lbl, anchor=south] at (n5.north) {pNext};
\node[ellipse, draw=red!80!black, thick, inner sep=3.5pt, fit=(d5)(n5), yshift=0.14cm] (b5) {};
\node[addr, anchor=south] at (b5.north) {0xa36};

% Mũi tên liên kết xanh lá giữa các Node
\draw[-{Stealth[length=2mm]}, thick, green!60!black] (n1.east) to[out=70, in=135, looseness=1.1] (b2.north west);
\draw[-{Stealth[length=2mm]}, thick, green!60!black] (n2.east) to[out=70, in=135, looseness=1.1] (b3.north west);
\draw[-{Stealth[length=2mm]}, thick, green!60!black] (n3.east) to[out=70, in=135, looseness=1.1] (b4.north west);
\draw[-{Stealth[length=2mm]}, thick, green!60!black] (n4.east) to[out=70, in=135, looseness=1.1] (b5.north west);

% pHead và pTail
\node[ptr, below=0.45cm of b1] {pHead};
\node[ptr, below=0.45cm of b5] {pTail};
\end{tikzpicture}
\caption{Kiến trúc danh sách liên kết đơn trong bộ nhớ.}
\end{figure}

% ==============================================================================
\section{Các Bước Xây Dựng Danh Sách Liên Kết Đơn}
% ==============================================================================

\subsection{Quy trình 4 bước xây dựng}

\begin{enumerate}[leftmargin=2.2em, itemsep=0.5em]
    \item \textbf{Bước \#1: Tạo cấu trúc Node}
    \begin{itemize}[itemsep=0.15em]
        \item Mỗi \texttt{Node} chứa dữ liệu (\texttt{int data}) và con trỏ trỏ tới Node kế tiếp (\texttt{Node* pNext}).
    \end{itemize}
\begin{lstlisting}
struct Node{
    int data;
    Node* pNext;
};
\end{lstlisting}

    \item \textbf{Bước \#2: Khởi tạo Node từ \texttt{value}} \textit{(Viết hàm chuyển \texttt{value} thành \texttt{Node})}
    \begin{itemize}[itemsep=0.15em]
        \item Cấp phát động một \texttt{Node} mới, gán \texttt{data = value}, con trỏ \texttt{pNext = NULL} và trả về địa chỉ \texttt{Node} vừa tạo (ví dụ: Node tại \texttt{0xb12} có \texttt{data = 10}, \texttt{pNext = NULL}).
    \end{itemize}
\begin{lstlisting}
Node* initNode(int value)
{
    Node* p = new Node;

    p->data = value;
    p->pNext = NULL;
    return p;
}
\end{lstlisting}

    \needspace{4.5cm}
    \item \textbf{Bước \#3: Tạo cấu trúc \texttt{List}}
    \begin{itemize}[itemsep=0.15em]
        \item Cấu trúc \texttt{List} quản lý danh sách thông qua hai con trỏ \texttt{pHead} và \texttt{pTail}.
    \end{itemize}
\begin{lstlisting}
struct List
{
    Node* pHead;
    Node* pTail;
};
\end{lstlisting}

    \item \textbf{Bước \#4: Khởi tạo \texttt{List}}
    \begin{itemize}[itemsep=0.15em]
        \item Khi danh sách mới tạo chưa có phần tử nào (danh sách rỗng), cả \texttt{pHead} và \texttt{pTail} đều trỏ về \texttt{NULL}.
    \end{itemize}
\begin{lstlisting}
void initList(List& l)
{
    l.pHead = NULL;
    l.pTail = NULL;
}
\end{lstlisting}
\end{enumerate}

\begin{figure}[H]
\centering
\begin{tikzpicture}[
    nbox/.style={rectangle, draw=blue!70!black, fill=white, minimum width=0.85cm, minimum height=0.5cm, font=\small, inner sep=2pt},
    pbox/.style={rectangle, draw=orange!80!black, fill=white, minimum width=0.85cm, minimum height=0.5cm, font=\footnotesize, inner sep=2pt},
    lbl/.style={font=\ttfamily\tiny, text=codeblue},
    addr/.style={font=\ttfamily\footnotesize, text=addrcolor}
]
% Bước 1: Khung Node mẫu
\node[font=\bfseries\small\color{codeblue}] at (0.45, 1.55) {Bước \#1: Cấu trúc Node};
\node[nbox] (s1d) at (0, 0) {\dots};
\node[pbox, right=0.05cm of s1d] (s1n) {\dots};
\node[lbl, anchor=south] at (s1d.north) {data};
\node[lbl, anchor=south] at (s1n.north) {pNext};

% Bước 2: Node 0xb12 (10 | NULL)
\begin{scope}[xshift=4.8cm]
\node[font=\bfseries\small\color{codeblue}] at (0.45, 1.55) {Bước \#2: Khởi tạo Node};
\node[nbox] (s2d) at (0, 0) {10};
\node[pbox, fill=yellow!80, right=0.05cm of s2d] (s2n) {NULL};
\node[lbl, anchor=south] at (s2d.north) {data};
\node[lbl, anchor=south] at (s2n.north) {pNext};
\node[ellipse, draw=red!80!black, thick, inner sep=3.5pt, fit=(s2d)(s2n), yshift=0.14cm] (s2b) {};
\node[addr, anchor=south] at (s2b.north) {0xb12};
\end{scope}

% Bước 3: Cấu trúc List
\begin{scope}[xshift=9.2cm]
\node[font=\bfseries\small\color{codeblue}] at (2.1, 1.55) {Bước \#3: Cấu trúc List};
\foreach \i in {0,1,2,3,4} {
    \node[ellipse, draw=orange!80!black, thick, fill=yellow!15, minimum width=0.75cm, minimum height=0.45cm] (ln\i) at (\i*1.05, 0.1) {};
}
\draw[thick, blue!70!black] (ln0) -- (ln1) -- (ln2) -- (ln3) -- (ln4);
\node[font=\bfseries\footnotesize, text=addrcolor, below=0.25cm of ln0] {pHead};
\node[font=\bfseries\footnotesize, text=addrcolor, below=0.25cm of ln4] {pTail};
\end{scope}
\end{tikzpicture}
\caption{Minh họa trực quan các bước cấu trúc Node và List.}
\end{figure}

\subsection{Khai báo -- Khởi tạo Node và List}
\begin{itemize}[itemsep=0.25em]
    \item \textbf{Khai báo danh sách và khởi tạo danh sách:}
\begin{lstlisting}
List l;
initList(l);
\end{lstlisting}
    Lúc này danh sách rỗng: \texttt{pHead => NULL}, \texttt{pTail => NULL}.

    \item \textbf{Khai báo Node và khởi tạo Node:}
\begin{lstlisting}
Node* p1 = initNode(10);
Node* p2 = initNode(79);
Node* p3 = initNode(39);
\end{lstlisting}
    Mỗi lần gọi \texttt{initNode}, một vùng nhớ động cho \texttt{Node} được cấp phát độc lập (ví dụ \texttt{p1} tại \texttt{0xb12} có \texttt{data = 10}, \texttt{pNext = NULL}).
\end{itemize}

\begin{figure}[H]
\centering
\begin{tikzpicture}[
    nbox/.style={rectangle, draw=blue!70!black, fill=white, minimum width=0.85cm, minimum height=0.5cm, font=\small, inner sep=2pt},
    pbox/.style={rectangle, draw=orange!80!black, fill=white, minimum width=0.85cm, minimum height=0.5cm, font=\footnotesize, inner sep=2pt},
    lbl/.style={font=\ttfamily\tiny, text=codeblue},
    addr/.style={font=\ttfamily\footnotesize, text=addrcolor}
]
% List rỗng
\node[font=\bfseries\small\color{codeblue}] at (0.45, 1.45) {Danh sách rỗng};
\node[nbox, text=red] (ed) at (0, 0.15) {\dots};
\node[pbox, fill=yellow!80, text=red, right=0.05cm of ed] (en) {\dots};
\node[lbl, anchor=south] at (ed.north) {data};
\node[lbl, anchor=south] at (en.north) {pNext};
\node[ellipse, draw=red!80!black, thick, inner sep=3.5pt, fit=(ed)(en), yshift=0.14cm] (eb) {};
\node[font=\ttfamily\footnotesize, text=codeblue, anchor=south] at (eb.north) {NULL};
\node[font=\bfseries\footnotesize, text=addrcolor, below=0.25cm of eb] (ep1) {pHead => \color{codeblue}NULL};
\node[font=\bfseries\footnotesize, text=addrcolor, below=0.05cm of ep1] {pTail => \color{codeblue}NULL};

% Node p1 khởi tạo
\begin{scope}[xshift=6.5cm]
\node[font=\bfseries\small\color{codeblue}] at (0.45, 1.45) {Node \texttt{p1 = initNode(10)}};
\node[nbox] (nd) at (0, 0.15) {10};
\node[pbox, fill=yellow!80, right=0.05cm of nd] (nn) {NULL};
\node[lbl, anchor=south] at (nd.north) {data};
\node[lbl, anchor=south] at (nn.north) {pNext};
\node[ellipse, draw=red!80!black, thick, inner sep=3.5pt, fit=(nd)(nn), yshift=0.14cm] (nb) {};
\node[addr, anchor=south] at (nb.north) {0xb12};
\end{scope}
\end{tikzpicture}
\caption{Trạng thái sau khi khởi tạo danh sách rỗng và khởi tạo một Node.}
\end{figure}

% ==============================================================================
\section{Các Thao Tác Trên Danh Sách Liên Kết Đơn}
% ==============================================================================

\subsection{Tổng quan quá trình thêm Node vào List}
Xét dãy giá trị cần đưa vào danh sách theo thứ tự: \texttt{10}, \texttt{79}, \texttt{39}, \texttt{50}, \texttt{26}. Với mỗi giá trị, ta tạo một Node mới bằng \texttt{Node* p = initNode(...);} rồi gọi \texttt{addHead(l, p)} (thêm vào đầu) hoặc \texttt{addTail(l, p)} (thêm vào cuối):

\begin{enumerate}[itemsep=0.25em]
    \item Giá trị \textbf{\texttt{10}}: Gọi \texttt{addTail(l, p);} hoặc \texttt{addHead(l, p);} (vì danh sách đang rỗng) $\Rightarrow$ \texttt{(10)}.
    \item Giá trị \textbf{\texttt{79}}: Gọi \texttt{addHead(l, p);} $\Rightarrow$ \texttt{(79) -- (10)}.
    \item Giá trị \textbf{\texttt{39}}: Gọi \texttt{addTail(l, p);} $\Rightarrow$ \texttt{(79) -- (10) -- (39)}.
    \item Giá trị \textbf{\texttt{50}}: Gọi \texttt{addHead(l, p);} $\Rightarrow$ \texttt{(50) -- (79) -- (10) -- (39)}.
    \item Giá trị \textbf{\texttt{26}}: Gọi \texttt{addTail(l, p);} $\Rightarrow$ \texttt{(50) -- (79) -- (10) -- (39) -- (26)}.
\end{enumerate}

\begin{figure}[H]
\centering
\begin{tikzpicture}[
    vnode/.style={ellipse, draw=orange!80!black, thick, fill=white, minimum width=0.95cm, minimum height=0.52cm, font=\small},
    cmd/.style={font=\ttfamily\footnotesize, fill=yellow!70, inner sep=2pt, rounded corners=2pt},
    lbl/.style={font=\bfseries\small, text=codeblue, anchor=east}
]
% Bước 1: 10
\node[lbl] at (-1.5, 0) {Bước 1 (10):};
\node[cmd] at (0.3, 0) {addTail / addHead};
\node[vnode] (s1_1) at (3.2, 0) {10};

% Bước 2: 79
\node[lbl] at (-1.5, -0.95) {Bước 2 (79):};
\node[cmd] at (0.3, -0.95) {addHead(l, p)};
\node[vnode] (s2_1) at (3.2, -0.95) {79};
\node[vnode] (s2_2) at (5.0, -0.95) {10};
\draw[thick, blue!70!black] (s2_1) -- (s2_2);

% Bước 3: 39
\node[lbl] at (-1.5, -1.9) {Bước 3 (39):};
\node[cmd] at (0.3, -1.9) {addTail(l, p)};
\node[vnode] (s3_1) at (3.2, -1.9) {79};
\node[vnode] (s3_2) at (5.0, -1.9) {10};
\node[vnode] (s3_3) at (6.8, -1.9) {39};
\draw[thick, blue!70!black] (s3_1) -- (s3_2) -- (s3_3);

% Bước 4: 50
\node[lbl] at (-1.5, -2.85) {Bước 4 (50):};
\node[cmd] at (0.3, -2.85) {addHead(l, p)};
\node[vnode] (s4_1) at (3.2, -2.85) {50};
\node[vnode] (s4_2) at (5.0, -2.85) {79};
\node[vnode] (s4_3) at (6.8, -2.85) {10};
\node[vnode] (s4_4) at (8.6, -2.85) {39};
\draw[thick, blue!70!black] (s4_1) -- (s4_2) -- (s4_3) -- (s4_4);

% Bước 5: 26
\node[lbl] at (-1.5, -3.8) {Bước 5 (26):};
\node[cmd] at (0.3, -3.8) {addTail(l, p)};
\node[vnode] (s5_1) at (3.2, -3.8) {50};
\node[vnode] (s5_2) at (5.0, -3.8) {79};
\node[vnode] (s5_3) at (6.8, -3.8) {10};
\node[vnode] (s5_4) at (8.6, -3.8) {39};
\node[vnode] (s5_5) at (10.4, -3.8) {26};
\draw[thick, blue!70!black] (s5_1) -- (s5_2) -- (s5_3) -- (s5_4) -- (s5_5);
\end{tikzpicture}
\caption{Diễn biến danh sách qua từng thao tác \texttt{addHead} và \texttt{addTail}.}
\end{figure}

% ------------------------------------------------------------------------------
\needspace{7.5cm}
\subsection{Thêm Node tại vị trí đầu List (\texttt{addHead})}
Hàm \texttt{addHead(List\& l, Node* p)} được cài đặt theo hai trường hợp: danh sách đang rỗng và danh sách đã có ít nhất một phần tử:
\begin{lstlisting}
void addHead(List& l, Node* p)
{
    if(l.pTail == NULL)
    {
        l.pHead = p;
        l.pTail = p;
    }
    else
    {
        p->pNext = l.pHead;
        l.pHead = p;
    }
}
\end{lstlisting}
\textit{(Ghi chú: Điều kiện \texttt{l.pTail == NULL} dùng để kiểm tra danh sách rỗng, tương đương với kiểm tra \texttt{l.pHead == NULL}.)}

\subsubsection*{Minh họa từng bước với dãy giá trị \texttt{10, 79, 39}:}

\noindent\textbf{Trường hợp 1: Danh sách đang rỗng (thêm giá trị \texttt{10})}
\begin{enumerate}[itemsep=0.15em]
    \item Ban đầu danh sách rỗng (\texttt{pHead => NULL}, \texttt{pTail => NULL}). Khởi tạo Node \texttt{p} tại địa chỉ \texttt{0xd56} bằng \texttt{initNode(10)} (\texttt{data = 10}, \texttt{pNext = NULL}).
    \item Do \texttt{l.pTail == NULL}, thực hiện nhánh \texttt{if}:
    \begin{center}
        \texttt{l.pHead = p; \quad l.pTail = p;}
    \end{center}
    \item Kết quả: Cả \texttt{pHead} và \texttt{pTail} cùng trỏ tới \texttt{0xd56}.
\end{enumerate}

\begin{figure}[H]
\centering
\begin{tikzpicture}[
    nbox/.style={rectangle, draw=blue!70!black, fill=white, minimum width=0.8cm, minimum height=0.48cm, font=\small, inner sep=2pt},
    pbox/.style={rectangle, draw=orange!80!black, fill=white, minimum width=0.8cm, minimum height=0.48cm, font=\footnotesize, inner sep=2pt},
    lbl/.style={font=\ttfamily\tiny, text=codeblue},
    addr/.style={font=\ttfamily\footnotesize, text=addrcolor}
]
% Trạng thái trước
\node[font=\bfseries\small\color{codeblue}] at (1.55, 2.85) {Trước khi thêm (\texttt{l.pTail == NULL})};
\node[nbox] (d0) at (0, 0.25) {\dots};
\node[pbox, fill=yellow!80, right=0.05cm of d0] (n0) {\dots};
\node[lbl, anchor=south] at (d0.north) {data};
\node[lbl, anchor=south] at (n0.north) {pNext};
\node[ellipse, draw=red!80!black, thick, inner sep=3.5pt, fit=(d0)(n0), yshift=0.14cm] (b0) {};
\node[font=\ttfamily\footnotesize, text=codeblue, anchor=south] at (b0.north) {NULL};
\node[font=\bfseries\footnotesize, text=addrcolor, below=0.25cm of b0] (t0) {pHead => \color{codeblue}NULL};
\node[font=\bfseries\footnotesize, text=addrcolor, below=0.02cm of t0] {pTail => \color{codeblue}NULL};

\node[nbox] (dp) at (2.8, 1.3) {10};
\node[pbox, fill=yellow!80, right=0.05cm of dp] (np) {NULL};
\node[lbl, anchor=south] at (dp.north) {data};
\node[lbl, anchor=south] at (np.north) {pNext};
\node[ellipse, draw=red!80!black, thick, inner sep=3.5pt, fit=(dp)(np), yshift=0.14cm] (bp) {};
\node[font=\ttfamily\footnotesize, text=codeblue, anchor=south] at (bp.north) {p => 0xd56};
\draw[-{Stealth[length=2.2mm]}, dashed, thick, blue] (bp.west) to[out=185, in=75] (b0.north east);

% Mũi tên chuyển trạng thái
\draw[-{Stealth[length=2.5mm]}, very thick, codeblue] (5.3, 0.75) -- (6.5, 0.75);

% Trạng thái sau
\begin{scope}[xshift=9.2cm]
\node[font=\bfseries\small\color{codeblue}] at (0.4, 2.85) {Sau khi thực hiện nhánh \texttt{if}};
\node[nbox] (d1) at (0, 0.5) {10};
\node[pbox, fill=yellow!80, right=0.05cm of d1] (n1) {NULL};
\node[lbl, anchor=south] at (d1.north) {data};
\node[lbl, anchor=south] at (n1.north) {pNext};
\node[ellipse, draw=red!80!black, thick, inner sep=3.5pt, fit=(d1)(n1), yshift=0.14cm] (b1) {};
\node[addr, anchor=south] at (b1.north) {0xd56};
\node[font=\bfseries\footnotesize, text=addrcolor, below=0.25cm of b1] (t1) {pHead => 0xd56};
\node[font=\bfseries\footnotesize, text=addrcolor, below=0.02cm of t1] {pTail => 0xd56};
\end{scope}
\end{tikzpicture}
\caption{Minh họa \texttt{addHead} khi danh sách rỗng (thêm \texttt{10} tại \texttt{0xd56}).}
\end{figure}

\noindent\textbf{Trường hợp 2: Danh sách đã có phần tử (tiếp tục thêm giá trị \texttt{79})}
\begin{enumerate}[itemsep=0.15em]
    \item Danh sách hiện tại có Node \texttt{0xd56} (\texttt{data = 10}, \texttt{pNext = NULL}) với \texttt{pHead => 0xd56}, \texttt{pTail => 0xd56}.
    \item Tạo Node mới \texttt{p = initNode(79);} tại địa chỉ \texttt{0xc95} (\texttt{data = 79}, \texttt{pNext = NULL}) rồi gọi \texttt{addHead(l, p);}.
    \item Vì \texttt{l.pTail != NULL}, thực hiện nhánh \texttt{else}:
    \begin{itemize}[itemsep=0.1em]
        \item \texttt{p->pNext = l.pHead;}: Gán địa chỉ \texttt{pHead} hiện tại (\texttt{0xd56}) vào trường \texttt{pNext} của Node \texttt{p} (\texttt{0xc95}).
        \item \texttt{l.pHead = p;}: Cập nhật \texttt{pHead} trỏ về Node mới \texttt{0xc95}.
    \end{itemize}
    \item Kết quả: \texttt{pHead => 0xc95}, \texttt{pTail => 0xd56}; Node \texttt{0xc95} nối tới Node \texttt{0xd56}.
\end{enumerate}

\begin{figure}[H]
\centering
\begin{tikzpicture}[
    nbox/.style={rectangle, draw=blue!70!black, fill=white, minimum width=0.78cm, minimum height=0.48cm, font=\small, inner sep=2pt},
    pbox/.style={rectangle, draw=orange!80!black, fill=white, minimum width=0.78cm, minimum height=0.48cm, font=\footnotesize, inner sep=2pt},
    lbl/.style={font=\ttfamily\tiny, text=codeblue},
    addr/.style={font=\ttfamily\footnotesize, text=addrcolor}
]
% Giai đoạn thực hiện p->pNext = l.pHead và l.pHead = p
\node[font=\bfseries\small\color{codeblue}] at (1.8, 1.75) {Thực hiện nhánh \texttt{else}};
\node[nbox] (d1) at (0, 0.35) {79};
\node[pbox, fill=yellow!80, right=0.05cm of d1] (n1) {NULL};
\node[lbl, anchor=south] at (d1.north) {data};
\node[lbl, anchor=south] at (n1.north) {pNext};
\node[ellipse, draw=red!80!black, thick, inner sep=3.5pt, fit=(d1)(n1), yshift=0.14cm] (b1) {};
\node[addr, anchor=south] at (b1.north) {0xc95};

\node[nbox] (d2) at (3.2, 0.35) {10};
\node[pbox, fill=yellow!80, right=0.05cm of d2] (n2) {NULL};
\node[lbl, anchor=south] at (d2.north) {data};
\node[lbl, anchor=south] at (n2.north) {pNext};
\node[ellipse, draw=red!80!black, thick, inner sep=3.5pt, fit=(d2)(n2), yshift=0.14cm] (b2) {};
\node[addr, anchor=south] at (b2.north) {0xd56};

\draw[-{Stealth[length=2mm]}, dashed, thick, green!60!black] (n1.east) to[out=70, in=135, looseness=1.1] (b2.north west);

\node[font=\bfseries\footnotesize, text=addrcolor, below=0.3cm of b1] (ph1) {pHead => ...?...};
\node[font=\bfseries\footnotesize, text=addrcolor, below=0.3cm of b2] (ph2) {pHead => 0xd56};
\node[font=\bfseries\footnotesize, text=addrcolor, below=0.02cm of ph2] {pTail => 0xd56};
\draw[-{Stealth[length=2mm]}, dotted, thick, green!60!black] (ph2.west) -- (ph1.east);

% Mũi tên chuyển trạng thái
\draw[-{Stealth[length=2.5mm]}, very thick, codeblue] (5.5, 0.4) -- (6.5, 0.4);

% Trạng thái kết quả
\begin{scope}[xshift=7.8cm]
\node[font=\bfseries\small\color{codeblue}] at (1.8, 1.75) {Trạng thái kết thúc};
\node[nbox] (rd1) at (0, 0.35) {79};
\node[pbox, fill=yellow!80, right=0.05cm of rd1] (rn1) {d56};
\node[lbl, anchor=south] at (rd1.north) {data};
\node[lbl, anchor=south] at (rn1.north) {pNext};
\node[ellipse, draw=red!80!black, thick, inner sep=3.5pt, fit=(rd1)(rn1), yshift=0.14cm] (rb1) {};
\node[addr, anchor=south] at (rb1.north) {0xc95};

\node[nbox] (rd2) at (3.2, 0.35) {10};
\node[pbox, fill=yellow!80, right=0.05cm of rd2] (rn2) {NULL};
\node[lbl, anchor=south] at (rd2.north) {data};
\node[lbl, anchor=south] at (rn2.north) {pNext};
\node[ellipse, draw=red!80!black, thick, inner sep=3.5pt, fit=(rd2)(rn2), yshift=0.14cm] (rb2) {};
\node[addr, anchor=south] at (rb2.north) {0xd56};

\draw[-{Stealth[length=2mm]}, thick, green!60!black] (rn1.east) to[out=70, in=135, looseness=1.1] (rb2.north west);

\node[font=\bfseries\footnotesize, text=addrcolor, below=0.3cm of rb1] {pHead => 0xc95};
\node[font=\bfseries\footnotesize, text=addrcolor, below=0.3cm of rb2] {pTail => 0xd56};
\end{scope}
\end{tikzpicture}
\caption{Minh họa \texttt{addHead} khi danh sách đã có phần tử (thêm \texttt{79} tại \texttt{0xc95}).}
\end{figure}

% ------------------------------------------------------------------------------
\subsection{Thêm Node tại vị trí cuối List (\texttt{addTail})}
Hàm \texttt{addTail(List\& l, Node* p)} được cài đặt tương ứng với hai trường hợp:
\begin{lstlisting}
void addTail(List& l, Node* p)
{
    if(l.pHead == NULL)
    {
        l.pHead = p;
        l.pTail = p;
    }
    else
    {
        l.pTail->pNext = p;
        l.pTail = p;
    }
}
\end{lstlisting}

\subsubsection*{Minh họa từng bước (tiếp tục thêm giá trị \texttt{39} vào cuối danh sách hiện có):}
\begin{enumerate}[itemsep=0.15em]
    \item Danh sách hiện tại có 2 Node: \texttt{0xc95 (79, d56)} $\rightarrow$ \texttt{0xd56 (10, NULL)}, với \texttt{pHead => 0xc95} và \texttt{pTail => 0xd56}.
    \item Khởi tạo Node mới \texttt{p = initNode(39);} tại địa chỉ \texttt{0xa36} (\texttt{data = 39}, \texttt{pNext = NULL}) và gọi \texttt{addTail(l, p);}.
    \item Vì \texttt{l.pHead != NULL}, thực hiện nhánh \texttt{else}:
    \begin{itemize}[itemsep=0.1em]
        \item \texttt{l.pTail->pNext = p;}: Cho trường \texttt{pNext} của Node cuối hiện tại (\texttt{0xd56}) lưu địa chỉ \texttt{a36} để liên kết tới Node mới \texttt{0xa36}.
        \item \texttt{l.pTail = p;}: Cập nhật \texttt{pTail} chuyển từ \texttt{0xd56} sang \texttt{0xa36}.
    \end{itemize}
    \item Kết quả: \texttt{pHead => 0xc95}, \texttt{pTail => 0xa36}; danh sách hoàn chỉnh gồm 3 Node: \texttt{0xc95 (79)} $\rightarrow$ \texttt{0xd56 (10)} $\rightarrow$ \texttt{0xa36 (39)}.
\end{enumerate}

\begin{figure}[H]
\centering
\begin{tikzpicture}[
    nbox/.style={rectangle, draw=blue!70!black, fill=white, minimum width=0.8cm, minimum height=0.46cm, font=\small, inner sep=2pt},
    pbox/.style={rectangle, draw=orange!80!black, fill=white, minimum width=0.8cm, minimum height=0.46cm, font=\footnotesize, inner sep=2pt},
    lbl/.style={font=\ttfamily\tiny, text=codeblue},
    addr/.style={font=\ttfamily\footnotesize, text=addrcolor}
]
% Giai đoạn thực hiện else (khối trên)
\node[font=\bfseries\small\color{codeblue}] at (3.4, 1.5) {Thực hiện \texttt{l.pTail->pNext = p; l.pTail = p;}};

\node[nbox] (d1) at (0, 0) {79};
\node[pbox, fill=yellow!80, right=0.05cm of d1] (n1) {d56};
\node[lbl, anchor=south] at (d1.north) {data};
\node[lbl, anchor=south] at (n1.north) {pNext};
\node[ellipse, draw=red!80!black, thick, inner sep=3.5pt, fit=(d1)(n1), yshift=0.13cm] (b1) {};
\node[addr, anchor=south] at (b1.north) {0xc95};

\node[nbox] (d2) at (3.0, 0) {10};
\node[pbox, fill=yellow!80, right=0.05cm of d2] (n2) {NULL};
\node[lbl, anchor=south] at (d2.north) {data};
\node[lbl, anchor=south] at (n2.north) {pNext};
\node[ellipse, draw=red!80!black, thick, inner sep=3.5pt, fit=(d2)(n2), yshift=0.13cm] (b2) {};
\node[addr, anchor=south] at (b2.north) {0xd56};

\node[nbox] (d3) at (6.0, 0) {39};
\node[pbox, fill=yellow!80, right=0.05cm of d3] (n3) {NULL};
\node[lbl, anchor=south] at (d3.north) {data};
\node[lbl, anchor=south] at (n3.north) {pNext};
\node[ellipse, draw=red!80!black, thick, inner sep=3.5pt, fit=(d3)(n3), yshift=0.13cm] (b3) {};
\node[addr, anchor=south] at (b3.north) {0xa36};

\draw[-{Stealth[length=2mm]}, thick, green!60!black] (n1.east) to[out=70, in=135, looseness=1.1] (b2.north west);
\draw[-{Stealth[length=2mm]}, dotted, thick, green!60!black] (n2.east) to[out=70, in=135, looseness=1.1] (b3.north west);

\node[font=\bfseries\footnotesize, text=addrcolor, below=0.3cm of b1] {pHead => 0xc95};
\node[font=\bfseries\footnotesize, text=addrcolor, below=0.3cm of b2] (pt1) {pTail => 0xd56};
\node[font=\bfseries\footnotesize, text=addrcolor, below=0.3cm of b3] (pt2) {pTail => ...?...};
\draw[-{Stealth[length=2mm]}, dotted, thick, blue] (pt1.east) -- (pt2.west);

% Trạng thái sau khi thêm (khối dưới)
\begin{scope}[yshift=-3.3cm]
\node[font=\bfseries\small\color{codeblue}] at (3.4, 1.5) {Trạng thái kết thúc};

\node[nbox] (rd1) at (0, 0) {79};
\node[pbox, fill=yellow!80, right=0.05cm of rd1] (rn1) {d56};
\node[lbl, anchor=south] at (rd1.north) {data};
\node[lbl, anchor=south] at (rn1.north) {pNext};
\node[ellipse, draw=red!80!black, thick, inner sep=3.5pt, fit=(rd1)(rn1), yshift=0.13cm] (rb1) {};
\node[addr, anchor=south] at (rb1.north) {0xc95};

\node[nbox] (rd2) at (3.0, 0) {10};
\node[pbox, fill=yellow!80, right=0.05cm of rd2] (rn2) {a36};
\node[lbl, anchor=south] at (rd2.north) {data};
\node[lbl, anchor=south] at (rn2.north) {pNext};
\node[ellipse, draw=red!80!black, thick, inner sep=3.5pt, fit=(rd2)(rn2), yshift=0.13cm] (rb2) {};
\node[addr, anchor=south] at (rb2.north) {0xd56};

\node[nbox] (rd3) at (6.0, 0) {39};
\node[pbox, fill=yellow!80, right=0.05cm of rd3] (rn3) {NULL};
\node[lbl, anchor=south] at (rd3.north) {data};
\node[lbl, anchor=south] at (rn3.north) {pNext};
\node[ellipse, draw=red!80!black, thick, inner sep=3.5pt, fit=(rd3)(rn3), yshift=0.13cm] (rb3) {};
\node[addr, anchor=south] at (rb3.north) {0xa36};

\draw[-{Stealth[length=2mm]}, thick, green!60!black] (rn1.east) to[out=70, in=135, looseness=1.1] (rb2.north west);
\draw[-{Stealth[length=2mm]}, thick, green!60!black] (rn2.east) to[out=70, in=135, looseness=1.1] (rb3.north west);

\node[font=\bfseries\footnotesize, text=addrcolor, below=0.3cm of rb1] {pHead => 0xc95};
\node[font=\bfseries\footnotesize, text=addrcolor, below=0.3cm of rb3] {pTail => 0xa36};
\end{scope}
\end{tikzpicture}
\caption{Minh họa \texttt{addTail} thêm Node \texttt{39} tại \texttt{0xa36} vào cuối danh sách.}
\end{figure}

% ------------------------------------------------------------------------------
\subsection{In danh sách (\texttt{printList} và \texttt{printWh})}
Để duyệt và xuất toàn bộ giá trị của danh sách (ví dụ danh sách \texttt{0xc95 (79)} $\rightarrow$ \texttt{0xd56 (10)} $\rightarrow$ \texttt{0xa36 (39)}), ta sử dụng một con trỏ \texttt{p} xuất phát từ \texttt{l.pHead} và di chuyển lần lượt theo \texttt{p = p->pNext} cho đến khi gặp \texttt{NULL}.

\begin{itemize}[itemsep=0.25em]
    \item \textbf{Cách 1: Duyệt bằng vòng lặp \texttt{for}}
\begin{lstlisting}
void printList(List l)
{
    for(Node* p = l.pHead; p != NULL; p = p->pNext)
    {
        cout << p->data << "\t";
    }
}
\end{lstlisting}

    \item \textbf{Cách 2: Duyệt bằng vòng lặp \texttt{while}}
\begin{lstlisting}
void printWh(List l)
{
    Node* p = l.pHead;
    while(p!=NULL)
    {
        cout << p->data << "\t";
        p = p->pNext;
    }
}
\end{lstlisting}
\end{itemize}

% ------------------------------------------------------------------------------
\needspace{8cm}
\subsection{Viết hàm \texttt{main} minh họa}
\begin{lstlisting}
int main()
{
    List l;
    initList(l);

    Node* p = initNode(10);

    addHead(l, p);
    addHead(l, initNode(79));
    addTail(l, initNode(39));

    cout << "In danh sach:\n";
    printList(l); //79  10  39

    cout << endl;
    printWh(l);
    return 0;
}
\end{lstlisting}

\noindent\textbf{Kết quả xuất ra màn hình:}
\begin{lstlisting}[numbers=none]
In danh sach:
79      10      39
79      10      39
\end{lstlisting}

% ------------------------------------------------------------------------------
\subsection{Tìm kiếm trên danh sách liên kết đơn}

\subsubsection{Tìm kiếm một giá trị có trong List hay không}
\begin{itemize}[itemsep=0.15em]
    \item \textbf{Yêu cầu:} Viết hàm tìm kiếm một giá trị có trong List hay không?
    \item \textbf{Phân tích:}
    \begin{itemize}
        \item \textit{Input:} List $\Rightarrow$ \texttt{List l}; Giá trị cần tìm $\Rightarrow$ \texttt{int value}.
        \item \textit{Output:} True/false $\Rightarrow$ \texttt{return bool}.
    \end{itemize}
    \item \textbf{Mở rộng:} Trả về địa chỉ (\texttt{Node*}).
\end{itemize}
\textbf{Khung cài đặt trên bài giảng:}
\begin{lstlisting}
bool timGiaTri(List l, int value)
{
    for(Node* p = l.pHead; p != NULL; p = p->pNext)
    {
        if(...........)
        {
            return true;
        }
    }

    return false;
}
\end{lstlisting}
\noindent\textit{Gợi ý hoàn thiện:} Thay \texttt{...........} bằng điều kiện \texttt{p->data == value}. Với phần mở rộng trả về địa chỉ, kiểu trả về của hàm là \texttt{Node*}; khi tìm thấy thì \texttt{return p;}, nếu duyệt hết danh sách không thấy thì \texttt{return NULL;}.

\subsubsection{Tìm giá trị Node kế cuối}
\begin{itemize}[itemsep=0.15em]
    \item \textbf{Yêu cầu:} Viết hàm tìm giá trị node kế cuối?
    \item \textbf{Phân tích:}
    \begin{itemize}
        \item \textit{Input:} List $\Rightarrow$ \texttt{List l}.
        \item \textit{Output:} Có hay không $\Rightarrow$ \texttt{return bool}; Giá trị tìm thấy $\Rightarrow$ giá trị.
    \end{itemize}
    \item \textbf{Mở rộng:} Trả về địa chỉ.
\end{itemize}

\begin{figure}[H]
\centering
\begin{tikzpicture}[
    n/.style={ellipse, draw=orange!80!black, thick, fill=yellow!20, minimum width=0.85cm, minimum height=0.45cm},
    ng/.style={ellipse, draw=orange!80!black, thick, fill=green!55!black, minimum width=0.85cm, minimum height=0.45cm},
    lbl/.style={font=\bfseries\footnotesize, text=addrcolor}
]
% Trường hợp 1: Nhiều node (>= 3 node)
\node[font=\bfseries\small, anchor=east] at (-0.9, 0) {Nhiều phần tử:};
\node[n] (a1) at (0, 0) {};
\node[n] (a2) at (1.15, 0) {};
\node[n] (a3) at (2.3, 0) {};
\node[ng] (a4) at (3.45, 0) {};
\node[n] (a5) at (4.6, 0) {};
\draw[thick, blue!70!black] (a1) -- (a2) -- (a3) -- (a4) -- (a5);
\node[lbl, below=0.15cm of a1] {pHead};
\node[lbl, below=0.15cm of a5] {pTail};
\node[font=\footnotesize\itshape, right=0.25cm of a5, text width=6.2cm] {(Node tô xanh là node kế cuối: \texttt{p->pNext == l.pTail})};

% Trường hợp 2: 2 node
\node[font=\bfseries\small, anchor=east] at (-0.9, -1.25) {2 phần tử:};
\node[ng] (b1) at (0, -1.25) {};
\node[n] (b2) at (1.15, -1.25) {};
\draw[thick, blue!70!black] (b1) -- (b2);
\node[lbl, below=0.15cm of b1] {pHead};
\node[lbl, below=0.15cm of b2] {pTail};
\node[font=\footnotesize\itshape, right=0.25cm of b2] {(Node kế cuối chính là \texttt{pHead})};

% Trường hợp 3: 1 node
\node[font=\bfseries\small, anchor=east] at (-0.9, -2.5) {1 phần tử:};
\node[n] (c1) at (0, -2.5) {};
\node[lbl, below=0.1cm of c1] (cp1) {pHead};
\node[lbl, below=0.01cm of cp1] {pTail};
\node[font=\footnotesize\itshape, right=0.25cm of c1] {(Không có node kế cuối)};
\end{tikzpicture}
\caption{Các trường hợp cần xét khi tìm Node kế cuối trong danh sách.}
\end{figure}

% ------------------------------------------------------------------------------
\subsection{Kỹ thuật đếm}
\begin{itemize}[itemsep=0.15em]
    \item \textbf{Yêu cầu:} Viết hàm đếm các giá trị chẵn có trong List?
    \item \textbf{Phân tích:}
    \begin{itemize}
        \item \textit{Input:} List $\Rightarrow$ \texttt{List l}.
        \item \textit{Output:} Số lượng phần tử chẵn $\Rightarrow$ \texttt{return int}.
    \end{itemize}
\end{itemize}
\noindent\textit{Hướng thực hiện:} Khởi tạo biến đếm \texttt{int dem = 0;}, duyệt từ \texttt{l.pHead} đến \texttt{NULL}; với mỗi Node \texttt{p}, nếu \texttt{p->data \% 2 == 0} thì tăng \texttt{dem} lên \texttt{1}, cuối cùng \texttt{return dem;}.

% ------------------------------------------------------------------------------
\needspace{10.5cm}
\subsection{Tìm giá trị lớn nhất -- nhỏ nhất}

\subsubsection{Tìm giá trị lớn nhất}
\begin{itemize}[itemsep=0.15em]
    \item \textbf{Yêu cầu:} Viết hàm tìm giá trị lớn nhất trong List?
    \item \textbf{Phân tích:}
    \begin{itemize}
        \item \textit{Input:} List $\Rightarrow$ \texttt{List l}.
        \item \textit{Output:} Giá trị lớn nhất $\Rightarrow$ \texttt{return} giá trị.
    \end{itemize}
\end{itemize}
\textbf{Khung cài đặt trên bài giảng:}
\begin{lstlisting}
int timGiaTriMax(List l)
{
    int m = (*@..giá trị đầu tiên...@*);
    for(Node* p = l.pHead; p != NULL; p = p->pNext)
    {
        if(............)
        {
            ............
        }
    }

    return m;
}
\end{lstlisting}
\noindent\textit{Gợi ý hoàn thiện:} Khởi tạo \texttt{int m = l.pHead->data;} (khi danh sách không rỗng); trong vòng lặp kiểm tra \texttt{if (p->data > m)} thì cập nhật \texttt{m = p->data;}.

\subsubsection{Tìm giá trị nhỏ nhất}
\begin{itemize}[itemsep=0.15em]
    \item \textbf{Yêu cầu:} Viết hàm tìm giá trị nhỏ nhất trong List?
    \item \textbf{Phân tích:}
    \begin{itemize}
        \item \textit{Input:} List $\Rightarrow$ \texttt{List l}.
        \item \textit{Output:} Giá trị nhỏ nhất $\Rightarrow$ \texttt{return} giá trị \textit{(ghi chú: slide gốc ghi nhầm chữ ``Giá trị lớn nhất'')}.
    \end{itemize}
    \item \textbf{Kết luận:} $\Rightarrow$ \textbf{Phương pháp tương tự tìm max} (giả sử \texttt{min} ban đầu là giá trị của \texttt{pHead}, sau đó duyệt qua từng Node đến \texttt{pTail} để so sánh và cập nhật).
\end{itemize}

\begin{figure}[H]
\centering
\begin{tikzpicture}[
    n/.style={ellipse, draw=orange!80!black, thick, fill=yellow!20, minimum width=0.85cm, minimum height=0.45cm},
    lbl/.style={font=\bfseries\footnotesize, text=addrcolor}
]
\node[n] (m1) at (0, 0) {};
\node[n] (m2) at (1.6, 0) {};
\node[n] (m3) at (3.2, 0) {};
\node[n] (m4) at (4.8, 0) {};
\node[n] (m5) at (6.4, 0) {};
\draw[thick, blue!70!black] (m1) -- (m2) -- (m3) -- (m4) -- (m5);

\node[lbl, below=0.15cm of m1] (ph) {pHead};
\node[lbl, below=0.15cm of m5] {pTail};
\node[font=\bfseries\footnotesize, text=codeblue, below=0.75cm of ph] (mn) {min};

\draw[-{Stealth[length=2mm]}, thick, blue!70!black] (ph.south) -- (mn.north);
\draw[-{Stealth[length=2mm]}, dashed, thick, green!60!black] (mn.east) -- (m2.south);
\draw[-{Stealth[length=2mm]}, dashed, thick, green!60!black] (mn.east) -- (m3.south);
\draw[-{Stealth[length=2mm]}, dashed, thick, green!60!black] (mn.east) -- (m4.south);
\draw[-{Stealth[length=2mm]}, dashed, thick, green!60!black] (mn.east) -- (m5.south);
\end{tikzpicture}
\caption{Minh họa phương pháp tìm giá trị nhỏ nhất (khởi tạo từ \texttt{pHead} và duyệt đến \texttt{pTail}).}
\end{figure}

% ------------------------------------------------------------------------------
\needspace{8.5cm}
\subsection{Tính tổng các giá trị trong List}
\begin{itemize}[itemsep=0.15em]
    \item \textbf{Yêu cầu:} Viết hàm tính tổng các giá trị trong List?
    \item \textbf{Phân tích:}
    \begin{itemize}
        \item \textit{Input:} List $\Rightarrow$ \texttt{List}.
        \item \textit{Output:} Tổng các giá trị $\Rightarrow$ \texttt{return} giá trị.
    \end{itemize}
\end{itemize}
\textbf{Khung cài đặt trên bài giảng:}
\begin{lstlisting}
int tinh(List l)
{
    int t = 0;
    for(Node* p = l.pHead; p != NULL; p = p->pNext)
    {
        t = (*@\color{codegreen}\itshape //công thức tính tổng@*)
    }

    return t;
}
\end{lstlisting}
\noindent\textit{Gợi ý hoàn thiện:} Công thức tính tổng là \texttt{t = t + p->data;}.

% ------------------------------------------------------------------------------
\needspace{7.5cm}
\subsection{Xóa Node tại vị trí đầu List (\texttt{removeHead})}
\begin{itemize}[itemsep=0.15em]
    \item \textbf{Yêu cầu:} Viết hàm xóa Node đầu tiên trong List?
    \item \textbf{Phân tích:}
    \begin{itemize}
        \item \textit{Input:} List $\Rightarrow$ \texttt{List}.
        \item \textit{Output:}
        \begin{itemize}
            \item List đã xóa $\Rightarrow$ \texttt{List\&}
            \item Kết quả xóa $\Rightarrow$ \texttt{bool}
            \item Giá trị đã xóa $\Rightarrow$ \textit{(trả về qua hàm hoặc tham chiếu)}
        \end{itemize}
    \end{itemize}
\end{itemize}
\textbf{Khung hàm trên bài giảng:}
\begin{lstlisting}
int removeHead(List& l)
{


    return value;
}
\end{lstlisting}

\subsubsection*{Minh họa từng bước thao tác xóa đầu:}
Xét danh sách ban đầu có 3 Node: \texttt{0xc95 (79, d56)} $\rightarrow$ \texttt{0xd56 (10, a36)} $\rightarrow$ \texttt{0xa36 (39, NULL)}, với \texttt{pHead => 0xc95}, \texttt{pTail => 0xa36}.
\begin{enumerate}[itemsep=0.2em]
    \item \textbf{Bước \fbox{\color{stepblue}\#1}:} Dùng con trỏ tạm \texttt{p} trỏ tới Node đầu tiên đang cần xóa (\texttt{p => 0xc95}, tương ứng \texttt{Node* p = l.pHead;} và lưu lại \texttt{int value = p->data;}).
    \item \textbf{Bước \fbox{\color{stepblue}\#2}:} Di chuyển \texttt{pHead} sang Node kế tiếp (\texttt{pHead} từ \texttt{0xc95} sang \texttt{0xd56}, tương ứng \texttt{l.pHead = l.pHead->pNext;}).
    \item \textbf{Bước \fbox{\color{stepblue}\#3}:} Hủy (thu hồi vùng nhớ) Node mà \texttt{p} đang trỏ tới (\texttt{delete p;}).
\end{enumerate}
\textit{(Lưu ý: Nếu danh sách trước khi xóa chỉ có 1 Node thì sau bước \#2, \texttt{l.pHead == NULL}, cần gán thêm \texttt{l.pTail = NULL;} để tránh lạc con trỏ.)}

\begin{figure}[H]
\centering
\begin{tikzpicture}[
    nbox/.style={rectangle, draw=blue!70!black, fill=white, minimum width=0.8cm, minimum height=0.46cm, font=\small, inner sep=2pt},
    pbox/.style={rectangle, draw=orange!80!black, fill=white, minimum width=0.8cm, minimum height=0.46cm, font=\footnotesize, inner sep=2pt},
    lbl/.style={font=\ttfamily\tiny, text=codeblue},
    addr/.style={font=\ttfamily\footnotesize, text=addrcolor},
    stag/.style={rectangle, draw=stepblue, thick, text=stepblue, font=\bfseries\footnotesize, inner sep=2pt},
    ptag/.style={rectangle, fill=green!45!yellow!75, text=addrcolor, font=\ttfamily\footnotesize\bfseries, inner sep=2.5pt}
]
% --- TRẠNG THÁI 1: Bước #1 và #2 (khối trên) ---
\node[font=\bfseries\small\color{codeblue}] at (3.4, 2.25) {Bước \#1 và Bước \#2};
\node[stag] (s1) at (-0.2, 1.55) {\#1};
\node[ptag, right=0.2cm of s1] {p => 0xc95};

\node[nbox] (a1) at (0, 0.25) {79};
\node[pbox, fill=yellow!80, right=0.05cm of a1] (b1) {d56};
\node[lbl, anchor=south] at (a1.north) {data};
\node[lbl, anchor=south] at (b1.north) {pNext};
\node[ellipse, draw=red!80!black, thick, inner sep=3.5pt, fit=(a1)(b1), yshift=0.13cm] (n1) {};
\node[addr, anchor=south] at (n1.north) {0xc95};

\node[nbox] (a2) at (3.0, 0.25) {10};
\node[pbox, fill=yellow!80, right=0.05cm of a2] (b2) {a36};
\node[lbl, anchor=south] at (a2.north) {data};
\node[lbl, anchor=south] at (b2.north) {pNext};
\node[ellipse, draw=red!80!black, thick, inner sep=3.5pt, fit=(a2)(b2), yshift=0.13cm] (n2) {};
\node[addr, anchor=south] at (n2.north) {0xd56};

\node[nbox] (a3) at (6.0, 0.25) {39};
\node[pbox, fill=yellow!80, right=0.05cm of a3] (b3) {NULL};
\node[lbl, anchor=south] at (a3.north) {data};
\node[lbl, anchor=south] at (b3.north) {pNext};
\node[ellipse, draw=red!80!black, thick, inner sep=3.5pt, fit=(a3)(b3), yshift=0.13cm] (n3) {};
\node[addr, anchor=south] at (n3.north) {0xa36};

\draw[-{Stealth[length=2mm]}, thick, green!60!black] (b1.east) to[out=70, in=135, looseness=1.1] (n2.north west);
\draw[-{Stealth[length=2mm]}, thick, green!60!black] (b2.east) to[out=70, in=135, looseness=1.1] (n3.north west);

\node[font=\bfseries\footnotesize, text=addrcolor, below=0.3cm of n1] (ph1) {pHead => 0xc95};
\node[font=\bfseries\footnotesize, text=addrcolor, fill=yellow!80, below=0.3cm of n2] (ph2) {pHead => 0xd56};
\node[font=\bfseries\footnotesize, text=addrcolor, below=0.3cm of n3] {pTail => 0xa36};
\draw[-{Stealth[length=2mm]}, dotted, thick, blue] (ph1.east) -- (ph2.west);
\node[stag, below=0.1cm of ph2] {\#2};

% --- TRẠNG THÁI 2: Bước #3 và Kết quả (khối dưới) ---
\begin{scope}[yshift=-4.1cm]
\node[font=\bfseries\small\color{codeblue}] at (3.4, 2.25) {Bước \#3 (Hủy \texttt{p}) và Kết quả};
\node[stag] (rs1) at (-0.2, 1.55) {\#1};
\node[ptag, right=0.2cm of rs1] {p => 0xc95};

\node[nbox] (ra1) at (0, 0.25) {79};
\node[pbox, fill=yellow!80, right=0.05cm of ra1] (rb1) {d56};
\node[lbl, anchor=south] at (ra1.north) {data};
\node[lbl, anchor=south] at (rb1.north) {pNext};
\node[ellipse, draw=red!80!black, dotted, thick, inner sep=3.5pt, fit=(ra1)(rb1), yshift=0.13cm] (rn1) {};
\node[addr, anchor=south] at (rn1.north) {0xc95};

\node[nbox] (ra2) at (3.0, 0.25) {10};
\node[pbox, fill=yellow!80, right=0.05cm of ra2] (rb2) {a36};
\node[lbl, anchor=south] at (ra2.north) {data};
\node[lbl, anchor=south] at (rb2.north) {pNext};
\node[ellipse, draw=red!80!black, thick, inner sep=3.5pt, fit=(ra2)(rb2), yshift=0.13cm] (rn2) {};
\node[addr, anchor=south] at (rn2.north) {0xd56};

\node[nbox] (ra3) at (6.0, 0.25) {39};
\node[pbox, fill=yellow!80, right=0.05cm of ra3] (rb3) {NULL};
\node[lbl, anchor=south] at (ra3.north) {data};
\node[lbl, anchor=south] at (rb3.north) {pNext};
\node[ellipse, draw=red!80!black, thick, inner sep=3.5pt, fit=(ra3)(rb3), yshift=0.13cm] (rn3) {};
\node[addr, anchor=south] at (rn3.north) {0xa36};

\draw[-{Stealth[length=2mm]}, dotted, thick, green!60!black] (rb1.east) to[out=70, in=135, looseness=1.1] (rn2.north west);
\node[stag] at (2.15, 1.15) {\#3};
\draw[-{Stealth[length=2mm]}, thick, green!60!black] (rb2.east) to[out=70, in=135, looseness=1.1] (rn3.north west);

\node[font=\bfseries\footnotesize, text=addrcolor, below=0.3cm of rn2] (rph2) {pHead => 0xd56};
\node[font=\bfseries\footnotesize, text=addrcolor, below=0.3cm of rn3] {pTail => 0xa36};
\node[stag, below=0.1cm of rph2] {\#2};
\end{scope}
\end{tikzpicture}
\caption{Minh họa từng bước xóa Node đầu danh sách (\texttt{removeHead}).}
\end{figure}

% ------------------------------------------------------------------------------
\needspace{7cm}
\subsection{Xóa Node tại vị trí cuối List (\texttt{removeTail})}
\begin{itemize}[itemsep=0.15em]
    \item \textbf{Yêu cầu:} Viết hàm xóa Node cuối trong List?
    \item \textbf{Phân tích:}
    \begin{itemize}
        \item \textit{Input:} List $\Rightarrow$ \texttt{List}.
        \item \textit{Output:}
        \begin{itemize}
            \item List đã xóa $\Rightarrow$ \texttt{List\&}
            \item Giá trị đã xóa $\Rightarrow$ \texttt{return}
        \end{itemize}
    \end{itemize}
\end{itemize}
\textbf{Khung hàm trên bài giảng:}
\begin{lstlisting}
int removeTail(List& l)
{


    return value;
}
\end{lstlisting}

\subsubsection*{Minh họa từng bước thao tác xóa cuối:}
Xét danh sách ban đầu có 4 Node: \texttt{0xc95 (79, d56)} $\rightarrow$ \texttt{0xd56 (10, a36)} $\rightarrow$ \texttt{0xa36 (39, f12)} $\rightarrow$ \texttt{0xf12 (26, NULL)}, với \texttt{pHead => 0xc95}, \texttt{pTail => 0xf12}.
\begin{enumerate}[itemsep=0.2em]
    \item \textbf{Bước \fbox{\color{stepblue}\#1}:} Tìm Node kế cuối (Node đứng ngay trước \texttt{pTail}) và giữ bằng con trỏ \texttt{q} (trong hình là \texttt{0xa36}, thỏa \texttt{q->pNext == l.pTail}).
    \item \textbf{Bước \fbox{\color{stepblue}\#2}:} Dùng con trỏ \texttt{p} trỏ tới Node cuối cùng cần xóa (\texttt{p => 0xf12}, tương ứng \texttt{Node* p = l.pTail;} và lưu \texttt{int value = p->data;}).
    \item \textbf{Bước \fbox{\color{stepblue}\#3}:} Cập nhật \texttt{pTail} trỏ về Node kế cuối \texttt{q} (\texttt{pTail} từ \texttt{0xf12} chuyển sang \texttt{0xa36}, tương ứng \texttt{l.pTail = q;}).
    \item \textbf{Bước \fbox{\color{stepblue}\#4}:} Ngắt liên kết tới Node cuối cũ (\texttt{l.pTail->pNext = NULL;}) và thu hồi vùng nhớ của Node \texttt{p} (\texttt{delete p;}).
\end{enumerate}

\begin{figure}[H]
\centering
\begin{tikzpicture}[
    nbox/.style={rectangle, draw=blue!70!black, fill=white, minimum width=0.72cm, minimum height=0.44cm, font=\small, inner sep=2pt},
    pbox/.style={rectangle, draw=orange!80!black, fill=white, minimum width=0.72cm, minimum height=0.44cm, font=\footnotesize, inner sep=2pt},
    lbl/.style={font=\ttfamily\tiny, text=codeblue},
    addr/.style={font=\ttfamily\footnotesize, text=addrcolor},
    stag/.style={rectangle, draw=stepblue, thick, text=stepblue, font=\bfseries\footnotesize, inner sep=2pt},
    ptag/.style={rectangle, fill=green!45!yellow!75, text=addrcolor, font=\ttfamily\footnotesize\bfseries, inner sep=2.5pt}
]
% Khối trên: Các bước #1, #2, #3, #4
\node[font=\bfseries\small\color{codeblue}] at (4.9, 2.65) {Diễn biến các bước \#1, \#2, \#3, \#4};

\node[stag] (st1) at (6.4, 2.05) {\#1};
\node[ptag, below=0.12cm of st1] {q => 0xa36};

\node[stag] (st2) at (9.4, 2.05) {\#2};
\node[ptag, below=0.12cm of st2] {p => 0xf12};

\node[nbox] (d1) at (0, 0) {79};
\node[pbox, fill=yellow!80, right=0.05cm of d1] (n1) {d56};
\node[lbl, anchor=south] at (d1.north) {data};
\node[lbl, anchor=south] at (n1.north) {pNext};
\node[ellipse, draw=red!80!black, thick, inner sep=3.2pt, fit=(d1)(n1), yshift=0.13cm] (b1) {};
\node[addr, anchor=south] at (b1.north) {0xc95};

\node[nbox] (d2) at (2.95, 0) {10};
\node[pbox, fill=yellow!80, right=0.05cm of d2] (n2) {a36};
\node[lbl, anchor=south] at (d2.north) {data};
\node[lbl, anchor=south] at (n2.north) {pNext};
\node[ellipse, draw=red!80!black, thick, inner sep=3.2pt, fit=(d2)(n2), yshift=0.13cm] (b2) {};
\node[addr, anchor=south] at (b2.north) {0xd56};

\node[nbox] (d3) at (5.9, 0) {39};
\node[pbox, fill=yellow!80, right=0.05cm of d3] (n3) {f12};
\node[lbl, anchor=south] at (d3.north) {data};
\node[lbl, anchor=south] at (n3.north) {pNext};
\node[ellipse, draw=red!80!black, thick, inner sep=3.2pt, fit=(d3)(n3), yshift=0.13cm] (b3) {};
\node[addr, anchor=south] at (b3.north) {0xa36};

\node[nbox] (d4) at (8.85, 0) {26};
\node[pbox, fill=yellow!80, right=0.05cm of d4] (n4) {NULL};
\node[lbl, anchor=south] at (d4.north) {data};
\node[lbl, anchor=south] at (n4.north) {pNext};
\node[ellipse, draw=red!80!black, dotted, thick, inner sep=3.2pt, fit=(d4)(n4), yshift=0.13cm] (b4) {};
\node[addr, anchor=south] at (b4.north) {0xf12};

\draw[-{Stealth[length=2mm]}, thick, green!60!black] (n1.east) to[out=70, in=135, looseness=1.1] (b2.north west);
\draw[-{Stealth[length=2mm]}, thick, green!60!black] (n2.east) to[out=70, in=135, looseness=1.1] (b3.north west);
\draw[-{Stealth[length=2mm]}, dotted, thick, green!60!black] (n3.east) to[out=70, in=135, looseness=1.1] (b4.north west);
\node[stag] at (7.95, 1.0) {\#4};

\node[font=\bfseries\footnotesize, text=addrcolor, below=0.3cm of b1] {pHead => 0xc95};
\node[font=\bfseries\footnotesize, text=addrcolor, below=0.3cm of b3] (pt1) {pTail => 0xa36};
\node[font=\bfseries\footnotesize, text=addrcolor, below=0.3cm of b4] (pt2) {pTail => 0xf12};
\draw[-{Stealth[length=2mm]}, dotted, thick, blue] (pt2.west) -- (pt1.east);
\node[stag, below=0.12cm of pt1] {\#3};

% Khối dưới: Trạng thái kết thúc
\begin{scope}[yshift=-4.05cm, xshift=1.45cm]
\node[font=\bfseries\small\color{codeblue}] at (3.4, 1.35) {Trạng thái sau khi xóa xong};

\node[nbox] (fd1) at (0, 0) {79};
\node[pbox, fill=yellow!80, right=0.05cm of fd1] (fn1) {d56};
\node[lbl, anchor=south] at (fd1.north) {data};
\node[lbl, anchor=south] at (fn1.north) {pNext};
\node[ellipse, draw=red!80!black, thick, inner sep=3.2pt, fit=(fd1)(fn1), yshift=0.13cm] (fb1) {};
\node[addr, anchor=south] at (fb1.north) {0xc95};

\node[nbox] (fd2) at (2.95, 0) {10};
\node[pbox, fill=yellow!80, right=0.05cm of fd2] (fn2) {a36};
\node[lbl, anchor=south] at (fd2.north) {data};
\node[lbl, anchor=south] at (fn2.north) {pNext};
\node[ellipse, draw=red!80!black, thick, inner sep=3.2pt, fit=(fd2)(fn2), yshift=0.13cm] (fb2) {};
\node[addr, anchor=south] at (fb2.north) {0xd56};

\node[nbox] (fd3) at (5.9, 0) {39};
\node[pbox, fill=yellow!80, right=0.05cm of fd3] (fn3) {NULL};
\node[lbl, anchor=south] at (fd3.north) {data};
\node[lbl, anchor=south] at (fn3.north) {pNext};
\node[ellipse, draw=red!80!black, thick, inner sep=3.2pt, fit=(fd3)(fn3), yshift=0.13cm] (fb3) {};
\node[addr, anchor=south] at (fb3.north) {0xa36};

\draw[-{Stealth[length=2mm]}, thick, green!60!black] (fn1.east) to[out=70, in=135, looseness=1.1] (fb2.north west);
\draw[-{Stealth[length=2mm]}, thick, green!60!black] (fn2.east) to[out=70, in=135, looseness=1.1] (fb3.north west);

\node[font=\bfseries\footnotesize, text=addrcolor, below=0.3cm of fb1] {pHead => 0xc95};
\node[font=\bfseries\footnotesize, text=addrcolor, below=0.3cm of fb3] {pTail => 0xa36};
\end{scope}
\end{tikzpicture}
\caption{Minh họa từng bước xóa Node cuối danh sách (\texttt{removeTail}).}
\end{figure}

% ==============================================================================
\section{Hệ Thống Bài Tập Ứng Dụng}
% ==============================================================================

\subsection{Bài Tập 1: Danh sách liên kết đơn quản lý các số nguyên}
\textbf{Đề bài:} Viết chương trình cài đặt danh sách liên kết đơn quản lý các số nguyên. Thực hiện theo yêu cầu sau:
\begin{enumerate}[itemsep=0.15em]
    \item Xây dựng cấu trúc node và cấu trúc danh sách liên kết đơn.
    \item Viết hàm khởi tạo danh sách liên kết đơn và hàm khởi tạo địa chỉ node từ một số nguyên.
    \item Viết hàm chèn node vào đầu danh sách.
    \item Viết hàm chèn node vào cuối danh sách.
    \item Viết hàm khởi tạo giá trị cho danh sách từ mảng một chiều.
    \item Viết hàm in giá trị danh sách đã nhập.
    \item Viết hàm in giá trị kèm địa chỉ của từng node trong danh sách.
    \item Viết hàm kiểm tra danh sách có rỗng hay không.
    \item Viết hàm trả về node thứ $n$ trong danh sách.
    \item Viết hàm tìm kiếm một node có giá trị $X$ trong danh sách.
    \item Viết hàm tìm kiếm \underline{các địa chỉ node} có giá trị $X$ trong danh sách.
    \item Viết hàm đếm số lượng các node có giá trị âm, dương.
    \item Viết hàm tìm node có giá trị âm lớn nhất, tìm node có giá trị lẻ nhỏ nhất trong danh sách.
    \item Viết hàm cập nhật giá trị cho một node có giá trị là $X$ bằng giá trị $Y$ mới.
    \item Viết hàm copy danh sách.
\end{enumerate}

\subsection{Bài Tập 2: Danh sách liên kết đơn quản lý sinh viên}
\textbf{Đề bài:} Viết chương trình cài đặt danh sách liên kết đơn quản lý sinh viên với các thao tác như sau:
\begin{enumerate}[itemsep=0.15em]
    \item Xây dựng cấu trúc dữ liệu có tên \textbf{\color{addrcolor}\texttt{SinhVien}} lưu trữ các thông tin sau: \texttt{maSV}, \texttt{hoTen}, \texttt{diemTB}, \texttt{xepLoai}\dots{} (hoặc lưu trữ thêm các thông tin khác nếu muốn).
    \item Xây dựng cấu trúc node để lưu trữ sinh viên và cấu trúc danh sách liên kết đơn quản lý sinh viên.
    \item Viết hàm khởi tạo danh sách liên kết đơn và hàm khởi tạo địa chỉ node từ kiểu \texttt{SinhVien}.
    \item Viết hàm thêm sinh viên vào đầu danh sách.
    \item Viết hàm thêm sinh viên vào cuối danh sách.
    \item Viết hàm khởi tạo dữ liệu cho danh sách sinh viên \textit{(Dữ liệu kiểm thử xem cuối bài)}.
    \item Viết hàm in danh sách sinh viên.
    \item Viết hàm cập nhật xếp loại học lực.
    \item Viết hàm tìm điểm trung bình cao nhất.
    \item Viết hàm tìm kiếm sinh viên theo mã sinh viên.
    \item Viết hàm tìm SV có trung bình cao nhất.
    \item Viết hàm tìm SV có điểm trung bình thấp nhất.
\end{enumerate}

\needspace{5.5cm}
\noindent\textbf{Dữ liệu Sinh Viên dùng để kiểm thử:}
\begin{lstlisting}[numbers=none]
{123, "Nguyen Van A", 9.1, ""};
{124, "Nguyen Van B", 8.8, ""};
{125, "Nguyen Van C", 9.1, ""};
{126, "Nguyen Van D", 2.1, ""};
{127, "Nguyen Van F", 9.7, ""};
{128, "Nguyen Van G", 4.1, ""};
{129, "Nguyen Van H", 8.3, ""};
{130, "Nguyen Van K", 7.9, ""};
\end{lstlisting}

\end{document}
